\chapter{Scoring and
Classification}
\label{chapter-7-scoring-and-classification}

\section{The Orchestration of
Intelligence}\label{the-orchestration-of-intelligence}

A primary challenge in engineering a multi-modal threat detection
platform is the synthesis of disparate analytical signals into a
unified, actionable verdict. PhishGuard achieves this through a
sophisticated scoring and classification engine that seamlessly bridges
the continuous probabilistic output of the XGBoost machine learning
model with discrete, rule-based heuristic indicators.

This architecture ensures that the system is not only highly accurate
but also highly transparent, providing Security Operations Center (SOC)
analysts and end-users with explicit, human-readable explanations for
every classification.

\section{The ML Predictor
Architecture}\label{the-ml-predictor-architecture}

The core inference engine is housed within the
\texttt{backend/ml/predictor.py} module. To guarantee high-throughput
and sub-millisecond inference times during production deployment, the
\texttt{PhishingPredictor} class is implemented utilizing a thread-safe
Singleton design pattern.

\subsection{Thread-Safe
Initialization}\label{thread-safe-initialization}

During the FastAPI application startup sequence, the
\texttt{PhishingPredictor} invokes its \texttt{\_\_new\_\_} method,
locking the execution thread (\texttt{threading.Lock()}) to instantiate
the XGBoost model exactly once. The model weights are loaded into memory
from \texttt{models/phishingModel.json}, alongside the strict
21-dimensional feature schema defined in \texttt{modelMetadata.json}.
Once initialized, the model state becomes immutable, allowing the
asynchronous web server to process concurrent inference requests across
multiple threads without encountering race conditions or requiring
continuous disk I/O operations.

\subsection{Inference Execution}\label{inference-execution}

When a URL analysis request is routed to the predictor, the
21-dimensional \texttt{FeatureSet} is cast into a strongly typed
\texttt{numpy.ndarray} (\texttt{dtype=np.float64}). The
\texttt{predict()} method invokes the XGBoost \texttt{predict\_proba}
function, returning a continuous probability bounded between 0.0 and
1.0, representing the mathematical likelihood that the input belongs to
the malicious class.

\section{The Phishing Scorer
Module}\label{the-phishing-scorer-module}

While the XGBoost probability serves as the primary ground truth for URL
inputs, it operates as an opaque numerical value. To resolve the
black-box interpretability crisis, the \texttt{backend/ml/scorer.py}
module introduces a parallel heuristic scoring engine that
reverse-engineers the threat context.

\subsection{Explanatory Heuristics}\label{explanatory-heuristics}

Regardless of whether the XGBoost model is active, the
\texttt{PhishingScorer} executes three discrete heuristic calculators to
generate the human-readable \texttt{factors} array that will be
presented to the user: 1. \textbf{\texttt{calculateUrlStructureScore}:}
Applies deterministic penalty weights to structural anomalies. For
example, a URL length exceeding 75 characters generates a scaling
penalty (\texttt{min((urlLength\ -\ 75)\ /\ 100,\ 0.3)}), while the
presence of raw IP addresses or Suspicious TLDs adds fixed penalties. 2.
\textbf{\texttt{calculateOsintScore}:} Evaluates the infrastructure
context. If \texttt{isNewlyRegistered} is true (domain age $<${}
30 days), a 0.35 penalty is recorded. Detection in third-party
blacklists adds up to a 0.5 penalty. 3.
\textbf{\texttt{calculateFeatureScore}:} Analyzes holistic combinations
that indicate sophisticated evasion, such as a newly registered domain
combined with the presence of suspicious keywords in the URI string
(\texttt{score\ +=\ 0.3}).

\subsection{Graceful Heuristic
Fallback}\label{graceful-heuristic-fallback}

In the event of a catastrophic failure where the compiled XGBoost model
cannot be loaded into memory, the \texttt{PhishingScorer} implements a
graceful degradation protocol. Instead of halting the application, it
relies entirely on the heuristic calculators, combining their raw
outputs using the strict constants defined in the
\texttt{ScoringWeights} dataclass: - \textbf{URL Structure:} 25\% weight
- \textbf{OSINT Derived:} 35\% weight - \textbf{Feature Based:} 40\%
weight

This fallback mechanism ensures the platform remains operational and
capable of identifying overt threats even during degraded system states.

\section{Risk Level Determination}\label{risk-level-determination}

The final continuous score (whether derived from the XGBoost probability
or the heuristic fallback) is mapped to discrete, actionable categories
using the strict boundaries defined in the \texttt{RISK\_THRESHOLDS}
dictionary:

\begin{itemize}
\item
  \textbf{Safe (Score $<${} 0.20):} The domain exhibits normal
  structural characteristics, extensive DNS history, and clean
  reputation metrics.
\item
  \textbf{Low Risk (Score $<${} 0.40):} Minor anomalies detected,
  often typical of newly registered legitimate businesses or
  non-standard tracking URLs.
\item
  \textbf{Medium Risk / Suspicious (Score $<${} 0.60):} The domain
  exhibits multiple characteristics commonly associated with phishing
  infrastructure, requiring explicit user verification.
\item
  \textbf{High Risk (Score $<${} 0.80):} Strong indications of
  malicious intent, such as obscured IPs, known bad TLDs, or missing MX
  records on a young domain.
\item
  \textbf{Critical (Score \(\ge\) 0.80):} Definitive malicious
  classification. The XGBoost model expresses supreme confidence in the
  threat, or the domain has been explicitly flagged by external threat
  intelligence APIs.
\end{itemize}

\subsection{Reason Prioritization}\label{reason-prioritization}

To prevent alert fatigue and ensure the user interface remains
uncluttered, the \texttt{\_prioritizeReasons} method deduplicates and
sorts the aggregated heuristic factors. It employs a keyword-based
sorting algorithm, ensuring that highly critical warnings (e.g.,
``malicious,'' ``ip address,'' ``newly registered'') are prioritized at
the top of the array, while lower-severity structural warnings are
pushed down or truncated. The final output is strictly limited to the
top 10 most relevant reasons.

